% Vietnamese translation for OI Public Library
% Source-File: IOI中国国家候选队论文/2008/Day2/4.张煜承《一类算法复合的方法》/一类算法复合的方法.ppt
\documentclass[11pt]{article}
\usepackage{oipl-vn}

\newcommand{\Oh}{\mathcal{O}}

\title{Bản trình chiếu: Một phương pháp kết hợp các thuật toán}
\author{Zhang Yucheng}
\date{2008}

\begin{document}
\maketitle

\origtitle{A method for algorithm composition: slides}
\sourcepath{IOI China 2008 / Day 2 / Zhang Yucheng / slides.ppt}

\section{Mạch trình bày}

PPT gốc chứa nhiều hình và chữ trong đối tượng đồ họa; phần trích được rõ
nhất là bài toán thao tác \texttt{B X}, \texttt{A Y}, các giới hạn
\(N=40000\), \(R=500000\), công thức \(x\bmod Y\), các khoảng
\([kY,(k+1)Y)\), và một số mốc độ phức tạp. Bản ghi chú này dựa trên các mốc
ấy và bài viết đi kèm để phục dựng mạch trình bày của slide.

Chủ đề của bản trình chiếu là kết hợp thuật toán: cùng một bài toán có thể có
nhiều thuật toán đều chưa đủ tốt trong trường hợp xấu nhất, nhưng mỗi thuật
toán lại mạnh trên một vùng dữ liệu. Nếu chia bài toán theo đúng đặc trưng
của các vùng ấy, ta có thể dùng từng thuật toán ở nơi nó hiệu quả nhất và thu
được độ phức tạp tổng thể tốt hơn.

\section{Bài toán thao tác trên tập}

Bài toán thứ nhất duy trì một tập \(S\), ban đầu rỗng. Có hai thao tác:
\begin{itemize}
  \item \texttt{B X}: chèn số \(X\) vào \(S\);
  \item \texttt{A Y}: hỏi số trong \(S\) có phần dư nhỏ nhất khi chia cho
  \(Y\).
\end{itemize}
Các thao tác được biết trước, nên được phép xử lý offline. Giới hạn của slide
là \(N\le 40000\) và \(1\le X,Y\le R=500000\).

Thuật toán trực tiếp duyệt mọi phần tử hiện có của \(S\) khi gặp truy vấn
\texttt{A Y}. Nó tốt nếu số lần chèn hoặc số lần hỏi rất ít, nhưng trường hợp
xấu nhất là \(\Oh(N^2)\). Một thuật toán khác duy trì sẵn đáp án cho từng giá
trị \(Y\) xuất hiện trong truy vấn: mỗi lần chèn \(X\), cập nhật đáp án của
mọi \(Y\) cần quan tâm. Cách này cũng có thể là \(\Oh(N^2)\), nhưng nếu số
giá trị \(Y\) cần duy trì nhỏ thì rất tốt.

\section{Thuật toán theo khoảng modulo}

Slide nhấn mạnh phép biến đổi
\[
  x=kY+r,\qquad 0\le r<Y.
\]
Muốn \(x\bmod Y\) nhỏ nhất tức là muốn \(r\) nhỏ nhất. Với một \(k\) cố định,
ta chỉ cần tìm phần tử nhỏ nhất của \(S\) trong khoảng \([kY,(k+1)Y)\). Vì
vậy truy vấn \texttt{A Y} có thể được trả lời bằng cách duyệt các khoảng độ
dài \(Y\), lấy phần tử nhỏ nhất trong mỗi khoảng, rồi chọn phần dư tốt nhất.

Để hỏi phần tử nhỏ nhất không nhỏ hơn một đầu mút \(a\), bài viết dùng một
cấu trúc offline dựa trên hợp nhất--tìm kiếm. Nếu xử lý thời gian ngược, thao
tác chèn ở chiều xuôi trở thành thao tác xóa hoặc gộp khoảng ở chiều ngược;
các khoảng có cùng ``phần tử kế tiếp'' được duy trì bằng DSU. Chi tiết cài
đặt không phải trọng tâm của slide, nên ở đây chỉ giữ lại ý chính: mỗi khoảng
\([kY,(k+1)Y)\) có thể được kiểm tra gần như hằng số.

Với cách này, một truy vấn \texttt{A Y} cần xét khoảng \(\Oh(R/Y)\). Khi
\(Y\) lớn, số khoảng ít và thuật toán rất nhanh; khi \(Y\) nhỏ, nó chậm. Đây
là nửa thứ nhất của phép kết hợp.

\section{Chia truy vấn theo ngưỡng}

Nửa thứ hai là thuật toán duy trì sẵn đáp án cho các \(Y\) nhỏ. Chọn một
ngưỡng \(K\). Với mọi \(Y\le K\), ta duy trì đáp án \(p(Y)\) khi xử lý thao
tác theo chiều xuôi. Mỗi lần chèn \(X\), cập nhật \(p(Y)\) cho mọi
\(1\le Y\le K\), mất \(\Oh(K)\); mỗi truy vấn nhỏ trả lời \(\Oh(1)\).

Với \(Y>K\), ta dùng thuật toán theo khoảng ở trên. Mỗi truy vấn lớn chỉ cần
\(\Oh(R/Y)\le \Oh(R/K)\) lần kiểm tra khoảng. Tổng thời gian là
\[
  \Oh\left(NK+\frac{NR}{K}\right).
\]
Cân bằng hai hạng tử bằng cách chọn \(K=\sqrt{R}\), ta được
\[
  \Oh(N\sqrt{R}).
\]
Đây là ví dụ rất sạch của kết hợp thuật toán: thuật toán duy trì tốt cho
\(Y\) nhỏ, thuật toán quét khoảng tốt cho \(Y\) lớn, và ngưỡng chia được chọn
sao cho hai chi phí cân bằng.

\section{Bài toán đếm hình chữ nhật}

Bài toán thứ hai cho \(N\) điểm trên mặt phẳng, hỏi có bao nhiêu hình chữ nhật
cạnh song song trục tọa độ mà bốn đỉnh đều thuộc tập điểm. Sau khi rời rạc
hóa tọa độ, mỗi điểm nằm trên một nút lưới. Một hình chữ nhật được xác định
bởi hai hàng và hai cột cùng có điểm.

Thuật toán dễ nghĩ: xét từng cặp hàng. Với hai hàng \(i,j\), đếm số cột mà
cả hai hàng đều có điểm; nếu số đó là \(s\), hai hàng tạo ra
\[
  \binom{s}{2}
\]
hình chữ nhật. Cách này là \(\Oh(N^2)\) trong trường hợp xấu nhất, nhưng rất
tốt khi số hàng ít hoặc khi ta chỉ cần xử lý một nhóm hàng nhỏ.

Để tận dụng đặc điểm ấy, đặt ngưỡng \(K\). Hàng có hơn \(K\) điểm gọi là hàng
loại A, còn lại là loại B. Vì tổng số điểm là \(N\), số hàng loại A nhỏ hơn
\(N/K\). Do đó các hình chữ nhật dùng hai hàng loại A, và cả phần dùng một
hàng loại A với một hàng loại B, có thể xử lý bằng cách xét các hàng loại A
và quét các điểm, với chi phí cỡ
\[
  \Oh\left(\frac{N^2}{K}\right).
\]

\section{Thuật toán riêng cho các hàng thưa}

Phần còn lại là các hình chữ nhật dùng hai hàng loại B. Mỗi hàng loại B có
không quá \(K\) điểm. Thay vì xét cặp hàng, ta xét cặp cột trên cùng một
hàng. Mỗi cặp điểm cùng hàng tạo ra một đoạn ngang \([a,b]\) theo chỉ số cột.
Nếu cùng đoạn \([a,b]\) xuất hiện trên \(s\) hàng khác nhau, nó đóng góp
\(\binom{s}{2}\) hình chữ nhật.

Vì mỗi hàng loại B có không quá \(K\) điểm, tổng số đoạn cần xét là
\(\Oh(KN)\). Có thể băm các đoạn để đếm số lần xuất hiện. Bài viết còn nêu
cách giảm bộ nhớ: cố định đầu mút phải \(b\), chỉ băm đầu mút trái của những
đoạn có cùng \(b\), nhờ đó vẫn giữ thời gian \(\Oh(KN)\) nhưng không cần lưu
toàn bộ \(\Oh(KN)\) đoạn cùng lúc.

Kết hợp hai phần, tổng thời gian là
\[
  \Oh\left(\frac{N^2}{K}+KN\right).
\]
Chọn \(K=\sqrt{N}\) cho độ phức tạp
\[
  \Oh(N^{1.5}).
\]
Ở đây cách chia không phải theo giá trị truy vấn rõ ràng như bài thứ nhất, mà
theo mật độ của từng hàng. Thuật toán cặp hàng mạnh ở nhóm hàng dày vì số
hàng dày ít; thuật toán băm đoạn mạnh ở nhóm hàng thưa vì mỗi hàng sinh ít
đoạn.

\section{Tổng kết}

Kết hợp thuật toán không có nghĩa là ghép tùy tiện nhiều kỹ thuật vào cùng
một chương trình. Cần xác định rõ:
\begin{enumerate}
  \item mỗi thuật toán mạnh trong điều kiện nào;
  \item bài toán có thể chia thành các phần song song theo điều kiện ấy hay
  không;
  \item chi phí ở ngưỡng chia có cân bằng được không;
  \item các phần sau khi chia có phủ hết bài toán và không đếm trùng hay bỏ
  sót hay không.
\end{enumerate}

Hai ví dụ của slide đều dùng các thuật toán riêng lẻ có trường hợp xấu nhất
chưa đủ tốt. Sau khi chia truy vấn theo \(Y\), hoặc chia hàng theo mật độ,
nhược điểm của thuật toán này rơi vào vùng mà thuật toán kia xử lý tốt hơn.
Đó là giá trị chính của phương pháp: thay vì tìm một thuật toán thống trị mọi
trường hợp, ta chọn tiêu chí phân rã làm cho các ưu điểm bổ sung cho nhau.

\end{document}
